\documentclass{article}
\usepackage{amsmath, amssymb, amsthm}
\usepackage{hyperref}
\usepackage{geometry}
\geometry{margin=1in}

\title{Erd\H{o}s Problem \#386}
\date{Last updated: 2025-08-31}
\author{Source: erdosproblems.com}

\begin{document}
\maketitle

\noindent\textbf{Status:} OPEN \quad \textbf{No prize}

\noindent\textbf{Tags:} number theory, binomial coefficients

\medskip
\noindent\textbf{Problem Statement:}

Let $2\leq k\leq n-2$. Can $\binom{n}{k}$ be the product of consecutive primes infinitely often? For example\[\binom{21}{2}=2\cdot 3\cdot 5\cdot 7.\]

\bigskip
\noindent\textbf{Remarks:}

Erd\H{o}s and Graham write that 'a proof that this cannot happen infinitely often for $\binom{n}{2}$ seems hopeless; probably this can never happen for $\binom{n}{k}$ if $3\leq k\leq n-3$.'

Weisenberg has provided four easy examples that show Erd\H{o}s and Graham were too optimistic here:\[\binom{7}{3}=5\cdot 7,\]\[\binom{10}{4}= 2\cdot 3\cdot 5\cdot 7,\]\[\binom{14}{4} = 7\cdot 11\cdot 13,\]and\[\binom{15}{6}=5\cdot 7\cdot 11\cdot 13.\]The known values of $n$ for which $\binom{n}{2}$ is the product of consecutive primes are $4,6,15,21,715$ (see A280992).

\bigskip
\noindent\small{Source: \url{https://www.erdosproblems.com/386}}
\end{document}
